
%pakages
\documentclass[12pt,a4paper]{article}
\usepackage[utf8]{inputenc}
\usepackage{graphicx}
\usepackage{geometry}
\usepackage{float}
\usepackage{xcolor}
\usepackage{fancyhdr}
\setlength{\headheight}{14.5pt}
\usepackage{array}
\usepackage{hyperref}
\usepackage{caption}
\usepackage{textcomp}
\usepackage{amsmath}
\usepackage{amssymb}
\usepackage{gensymb}
\geometry{
    top=1in,
    bottom=1in,
    left=1in,
    right=1in
}
\pagestyle{fancy}
\fancyhf{}

\lhead{Project Progress Log}
\rhead{2026}
\cfoot{\thepage}

%Introduction

\begin{document}
\begin{titlepage}
    \centering

    \vspace*{2cm}

    {\Huge \textbf{PROJECT PROGRESS LOG}}\\[1cm]

    {\LARGE \textbf{UGV-1}}\\[1.5cm]

     \textbf{Team Members}\\[0.5cm]

    {\large
    G.N.S.S.Rithvik\\
    R.Srikar Varshit
    }\\[1.5cm]

    \textbf{Roll Numbers} \\[0.5cm]
    
    {\large
    EE26BTECH11211\\
    EE26BTECH11223
    }\\[1.5cm]
   
    \vfill

    {\large \today}

\end{titlepage}

\section{Assembling:}

\begin{figure}[H]
    \centering
    \includegraphics[width=0.65\linewidth]{WhatsApp Image 2026-11-02 at 6.42.22 PM.jpeg}
    \includegraphics[width=0.75\linewidth]{WhatsApp Image 2026-11-02 at 6.42.21 PM (1).jpeg}
    \caption{Initial connection}
    \label{fig:placeholder}
\end{figure}

\subsection{Workdone:}
\textbf{Initially we connected the arduino to the Motor Driver on the breadboard. Then we took some jumpers and we connected the ESP-32 to the Motor Driver and arduino and left it hanging in mid air.}

\begin{figure}[H]
    \centering
    \includegraphics[width=0.75\linewidth]{WhatsApp Image 2026-11-02 at 6.42.21 PM.jpeg}
    \caption{Initial problem faced}
    \label{fig:placeholder}
\end{figure}

\subsection{Observation:}
\textbf{After completing the connection we observed that the ESP-32 wouldn't glow. Then we thought there there might be a loose connection (or) there might be some non working jumber wires. }

\subsection{Workdone:}
\textbf{Then we removed the whole set up shown in figure 2.Then we thought of giving direct connection to ESP-32 from arduino and Motor Driver on the breadboard, not using jumper wires in between. }

\begin{figure}[H]
    \centering
    \includegraphics[angle=90,width=0.85\linewidth]{WhatsApp Image 2026-11-02 at 7.53.56 PM.jpeg}
    \caption{Assembling the ESP-32 and the motor driver on the breadboard}
    \label{fig:placeholder}
\end{figure}

\subsection{Workdone:}
\textbf{We assembled the ESP-32 and the motor driver on the breadboard as shown in Figure 3. The red wire in the figure 3 is the wire connecting to the GND of the ESP-32. }

\begin{figure}[H]
    \centering
    \includegraphics[width=0.75\linewidth]{WhatsApp Image 2026-11-02 at 7.53.56 PM (1).jpeg}
    \caption{Using Arduino and connecting it to the breadboard}
    \label{fig:placeholder}
\end{figure}

\subsection{Workdone:}
\textbf{Then we connected the arduino to the breadboard. We are using the arduino as a voltage stabilizer. It receives the high voltage by the batteries as power input and gives as 5V output to the ESP-32 as power.}

\subsection{Workdone:}
\textbf{We powered the ESP-32 with arduino as shown in Figure 4.}

\begin{figure}[H]
    \centering
    \includegraphics[width=0.65\linewidth]{WhatsApp Image 2026-11-02 at 7.53.56 PM (2).jpeg}
    \caption{Complete connection of the circuit}
    \label{fig:placeholder}
\end{figure}

\subsection{Workdone:}
\textbf{After giving power from the arduino to ESP-32 and the Motor Driver. }

\subsection{Observation:}
\textbf{We observed that there is short circuit occuring when given power.It got short circuited when we gave direct connection between the batteries and ESP-32,Motor Driver. We observed that the led on the ESP-32 is glowing continously, Unlike when we gave it a direct connection.}




\end{document}
